% ============ 20 光滑基变换定理 (原书 pp.127-129) ============
\chapter{光滑基变换定理}\label{sec:20}
\index{光滑基变换定理}
\index{紧合-光滑基变换定理}
\index{特殊化}
\index{逆特殊化映射}
\index{提升到特征零}
\index{初等纤维化}

这只是简要总结.

\subsection*{拓扑中的紧合-光滑基变换定理}

设 $\varphi \colon X \to S$ 是流形的紧合映射.若 $\varphi$ 光滑，即切空间上的映射 $d\varphi$ 在 $X$ 的所有点处都满，则 $\varphi$ 是纤维化，由此推出对 $X$ 上的任何局部常值层 $F$，$R^{r}\varphi_{*}F$ 局部常值.

\subsection*{光滑基变换定理}

本小节余下部分中，我们只考虑代数闭域 $k$ 上的簇.

考虑图：
\[
\begin{tikzcd}
X \arrow[d, "\varphi"'] & X' = X \times_{S} T \arrow[l, "f'"'] \arrow[d, "\varphi'"] \\
S & T \arrow[l, "f"']
\end{tikzcd}
\]
对 $X$ 上的任何层 $F$，有 $T$ 上层 的典范态射
\[
f^{*}(R^{r}\varphi_{*}F) \to R^{r}\varphi'_*(f'^{*}F),
\]
称为基变换态射(或映射)---见定理~\ref{thm:17.3} 的证明.

\begin{theorem}\label{thm:20.1}[光滑基变换]
若 $f \colon T \to S$ 光滑，则对挠与 $k$ 的特征互素的所有可构造层 $F$，基变换态射都是同构.
\end{theorem}

\subsection*{紧合-光滑基变换定理}

回忆紧合基变换定理的有限性部分说：若 $\varphi \colon X \to S$ 紧合且 $F$ 可构造，则 $R^{r}\varphi_{*}F$ 可构造.

\begin{theorem}\label{thm:20.2}[紧合-光滑基变换]
若 $\varphi \colon X \to S$ 紧合且光滑，$F$ 是带有限茎的局部常值层，则只要 $F$ 中的挠与 $k$ 的特征互素，$R^{r}\varphi_{*}F$ 就是带有限茎的局部常值层.
\end{theorem}

设 $x_0$ 与 $x_1$ 是看作概形的 $X$ 的点(因此 $x_0$ 与 $x_1$ 不必闭).我们说 $x_0$ 是 $x_1$ 的特殊化，如果它包含在 $x_1$ 的闭包中.则 $X$ 的每个含 $x_0$ 的 Zariski 开子集也含 $x_1$.选取几何点 $\bar{x}_0 \to x_0 \hookrightarrow X$ 与 $\bar{x}_1 \to x_1 \hookrightarrow X$.则 $\bar{x}_0$ 的平展邻域 $(U, u)$ 可以赋予 $\bar{x}_1$ 的平展邻域的结构($U \to X$ 的像含 $x_1$，因此可以选取提升 $\bar{x}_1 \to X$ 的态射 $\bar{x}_1 \to U$).一旦对 $\bar{x}_0$ 的每个平展邻域相容地做到这一点，就得到 $X_{\mathrm{et}}$ 上每个层 $F$ 的逆特殊化映射
\[
F_{\bar{x}_0} \to F_{\bar{x}_1}.
\]
例如，取 $F = \mathbb{G}_{a}$，得到严格局部环上的映射
\[
\mathcal{O}_{X, \bar{x}_0} \to \mathcal{O}_{X, \bar{x}_1}.
\]
逆特殊化映射的定义涉及选择，但一旦固定映射 $\mathcal{O}_{X, \bar{x}_0} \to \mathcal{O}_{X, \bar{x}_1}$，映射 $F_{\bar{x}_0} \to F_{\bar{x}_1}$ 就关于 $F$ 成为典范的.换言之，不存在唯一的典范映射 $F_{\bar{x}_0} \to F_{\bar{x}_1}$，而是存在一族特殊的映射，其成员称为逆特殊化映射.若族中一个成员是同构，则它们都是.

\begin{proposition}\label{prop:20.3}
$X_{\mathrm{et}}$ 上的可构造层 $F$ 局部常值，当且仅当逆特殊化映射 $F_{\bar{x}_0} \to F_{\bar{x}_1}$ 都是同构.
\end{proposition}

把它与如下拓扑情形比较.设 $F$ 是挖去原点的圆盘 $U$ 上的局部常值层，$j \colon U \hookrightarrow X$ 是 $U$ 到整个圆盘的包含.层 $F$ 对应于带 $\pi_1(U, u)$ 作用($u \in U$)的模 $M$.$j_{*}F$ 在 $u$ 处的茎是 $M$，在 $o$ 处的茎是 $M^{\pi_1(U,u)}$.于是，特殊茎 $(j_{*}F)_o$ 同构于一般茎 $(j_{*}F)_u$(通过逆特殊化映射)，当且仅当 $\pi_1$ 平凡作用于 $M$，这意味着 $j_{*}F$ 常值.

现在可以如下重述紧合-光滑基变换定理.

\begin{theorem}\label{thm:20.4}
设 $\varphi \colon X \to S$ 紧合且光滑，$F$ 是 $X$ 上带有限茎的局部常值层，其挠与 $k$ 的特征互素.对满足 $s_0$ 是 $s_1$ 的特殊化的一对几何点 $\bar{s}_0$ 与 $\bar{s}_1$，逆特殊化映射
\[
H^{r}(X_{\bar{s}_0}, F) \to H^{r}(X_{\bar{s}_1}, F)
\]
是同构.这里 $X_{\bar{s}} = X \times_{S} \bar{s}$，即 $\varphi$ 在 $\bar{s}$ 上的几何纤维.
\end{theorem}

\subsection*{应用}

上述结果对概形也成立，而这里有一个引人注目的应用.设 $X$ 是特征 $p \neq 0$ 的代数闭域 $k$ 上的完备非奇异簇.我们说 $X$ 可以提升到特征零，如果
\begin{enumerate}
  \item 存在离散赋值环 $R$，其分式域 $K$ 特征为零，剩余域为 $k$；且
  \item 存在在 $S = \operatorname{Spec} R$ 上紧合且光滑，特殊纤维为 $X$ 的概形 $\varphi \colon X \to S$.
\end{enumerate}
例如，$\mathbb{P}^{n}$ 的子簇 $X$ 可以提升到特征零，如果存在齐次多项式 $f_i(T_0, \ldots, T_n) \in R[T_0, \ldots, T_n]$，使得
\begin{enumerate}
  \item 模 $\mathfrak{m}_{R}$，$f_i$ 生成 $X$ 在 $\mathbb{P}^{n}$ 中的嵌入的齐次理想；
  \item 看作 $K[T_0, \ldots, T_n]$ 中多项式时，$f_i$ 定义与 $X$ 有相同维数的 $K$ 上的簇 $X_1$.
\end{enumerate}
显然，$\mathbb{P}^{n}$ 中的任何非奇异超曲面都可以提升---只需把定义 $X$ 的单个多项式从 $k[T_0, \ldots, T_n]$ 提升到 $R[T_0, \ldots, T_n]$.类似地，$\mathbb{P}^{n}$ 中的任何光滑完全交都可以提升到特征零.曲线与阿贝尔簇可以提升到特征零，但除此之外所知甚少.当然，许多簇不能提升.问题是这样的：假定 $X$ 在 $\mathbb{P}^{n}$ 中有余维数 $r$，但它的齐次理想 $\mathcal{I}(X)$ 需要 $s > r$ 个生成元；当 $\mathcal{I}(X)$ 的生成元提升到 $R[T_0, \ldots, T_n]$ 时，一般地它们将在特征零中定义维数小于 $X$ 的簇 $X_1$---一般只能说
\[
n - s \leq \dim X_1 \leq \dim X = n - r.
\]

\begin{theorem}\label{thm:20.5}
假定特征 $p \neq 0$ 的代数闭域 $k$ 上的簇 $X_0$ 可以提升到特征零域 $K$ 上的簇 $X_1$.对任何有限 Abel 群 $\Lambda$，
\[
H^{r}(X_0, \Lambda) \cong H^{r}(X_{1, K^{\mathrm{al}}}, \Lambda).
\]
特别地，$X_0$ 的 Betti 数与 $X_1$ 的 Betti 数相等.在下一节中，我们将证明 $X_1$ 的上同调群等于拓扑空间 $X_1(\mathbb{C})$ 的上同调群(假定 $K$ 可以嵌入 $\mathbb{C}$).
\end{theorem}
